\pdfbookmark[0]{Anzeigen. Kries: Mathematische Geographie. 1814.}{anzeigen-kries-1814}
\begin{center}
\rule{.85\textwidth}{.3pt}\par
Göttingische gelehrte Anzeigen.\quad 13. Juni 1814.\par
\rule{.85\textwidth}{.3pt}
\end{center}
\hspace*{\parindent}\textit{Lehrbuch der mathematischen Geographie von} \textsc{Friedrich Kries}, Professor am Gymnasium in Gotha. Mit sieben Kupfertafeln. 236 Seiten im Oktavformat. Leipzig, bei \textsc{G. J. Göschen}.

Mit diesem Lehrbuch für eine Wissenschaft, die für jeden gebildeten Menschen in so vieler Hinsicht von Interesse ist, wollte der Verfasser einen geeigneten Mittelweg finden. Auf der einen Seite stehen die dürftigen und oberflächlichen Abrisse, die den Lehrbüchern der politischen Geografie gewöhnlich vorangestellt werden. Sie zählen nur die wichtigsten Ergebnisse auf, ohne sie mathematisch zu begründen oder zu erläutern. Auf der anderen Seite stehen die umfangreicheren Werke, die vertiefte und weniger verbreitete Kenntnisse der höheren Mathematik voraussetzen. Von einem solchen Werk erwartet man nur eine zweckmäßige Auswahl aus dem Bekannten, die geordnet, gründlich und klar dargestellt wird. Neue Erkenntnisse, die die Wissenschaft selbst weiterbringen, erwartet man dabei nicht. Dieses Ziel hat der Verfasser tatsächlich erreicht. In zehn Abschnitten behandelt er die Gestalt des Erdkörpers im Allgemeinen, die mathematische Einteilung der Erdkugel und ihre Größe, ihre Drehung um ihre Achse und die damit verbundenen Erscheinungen sowie die Mittel, die geografische Breite eines Ortes zu bestimmen und eine Meridianlinie zu ziehen. Er behandelt die Bewegung der Erde um die Sonne, die Einteilung der Himmels- und der Erdkugel in Bezug auf diese Bewegung und die Erscheinungen, die sie auf der Erde hervorruft, sowie die Zeitbestimmung und die Mittel zur Bestimmung der geografischen Länge. Die letzten Abschnitte behandeln die sphäroidische Gestalt der Erde, die Herstellung von Erdgloben und Landkarten und den Gebrauch eines Erdglobus zur Lösung mathematisch-geografischer Aufgaben. Wir können dieser Anordnung und Auswahl im Allgemeinen nur zustimmen, obwohl unserer Ansicht nach hier und da noch einige Themen hätten aufgenommen werden können, die unerwähnt bleiben. Andere, etwa die verschiedenen Arten der Kartenprojektion, hätten dafür vielleicht etwas kürzer behandelt werden können. So hätten wir uns zum Beispiel eine Anleitung gewünscht, wie man den Flächeninhalt einzelner Länder berechnet,
